% One Fourier--Motzkin step, on the wire carrying the eliminated variable t.
% Part 1: a single row a_i t + u_i >= 0 is brought to the form t >= m_i
%         (i in P) or t <= m_j (j in N).
% Part 2: the bounds are merged pairwise by the generalised spider.
% Part 3: the result is relinearised and rescaled to the new cone M'.
\begin{center}
    \begin{tikzpicture}[axpic]
        \pic[val=a_i] (g1) at (0.500,1.000) {scalar};
        \pic (g2) at (2.000,0.500) {multiplication};
        \pic (g3) at (3.250,0.500) {ge};
        \pic (g4) at (4.500,0.500) {cozero};
        \coordinate (id_in_1)  at (0.000,0.000);
        \coordinate (id_out_1) at (1.250,0.000);
        \draw (id_in_1) -- (id_out_1);
        \draw (g1-out-0) to[out=0,in=180] (g2-in-0);
        \draw (id_out_1) to[out=0,in=180] (g2-in-1);
        \draw (g2-out-0) to[out=0,in=180] (g3-in-0);
        \draw (g3-out-0) to[out=0,in=180] (g4-in-0);
        \node[font=\scriptsize] at (0.100,1.200) {$t$};
        \node[font=\scriptsize] at (0.150,-0.200) {$u_i$};
    \end{tikzpicture}
    \;\begin{tikzpicture}[axpic]\node{$\overset{(\text{cap form})}{=}$};\end{tikzpicture}\;
    \begin{tikzpicture}[axpic]
        \pic[val=a_i] (g1) at (0.500,1.000) {scalar};
        \pic (g2) at (1.750,1.000) {ge};
        \pic (g3) at (1.750,0.000) {antipode};
        \coordinate (id_in_1)  at (0.000,0.000);
        \draw (id_in_1) -- (g3-in-0);
        \draw (g1-out-0) to[out=0,in=180] (g2-in-0);
        \draw (g2-out-0) -- (2.850,1.000) arc (90:-90:0.500) -- (g3-out-0);
        \node[font=\scriptsize] at (0.100,1.200) {$t$};
        \node[font=\scriptsize] at (0.150,-0.200) {$u_i$};
    \end{tikzpicture} \\
    \vspace{4mm}

    \;\begin{tikzpicture}[axpic]\node{$\overset{(P5),(P6),(\text{slide})}{=}$};\end{tikzpicture}\;
    \begin{tikzpicture}[axpic]
        \pic (g1) at (0.500,1.000) {ge};
        \pic[val=a_i^{-1}] (g2) at (0.900,0.000) {scalar};
        \pic (g3) at (2.150,0.000) {antipode};
        \coordinate (id_in_1) at (0.000,0.000);
        \draw (id_in_1) -- (g2-in-0);
        \draw (g2-out-0) to[out=0,in=180] (g3-in-0);
        \draw (g1-out-0) -- (3.050,1.000) arc (90:-90:0.500) -- (g3-out-0);
        \node[font=\scriptsize] at (0.100,1.200) {$t$};
        \node[font=\scriptsize] at (0.150,-0.250) {$u_i$};
        \node[font=\scriptsize] at (1.500,-0.750) {$i\in P$};
    \end{tikzpicture}
    \qquad
    \begin{tikzpicture}[axpic]
        \pic (g1) at (0.500,1.000) {coge};
        \pic[val=a_j^{-1}] (g2) at (0.900,0.000) {scalar};
        \pic (g3) at (2.150,0.000) {antipode};
        \coordinate (id_in_1) at (0.000,0.000);
        \draw (id_in_1) -- (g2-in-0);
        \draw (g2-out-0) to[out=0,in=180] (g3-in-0);
        \draw (g1-out-0) -- (3.050,1.000) arc (90:-90:0.500) -- (g3-out-0);
        \node[font=\scriptsize] at (0.100,1.200) {$t$};
        \node[font=\scriptsize] at (0.150,-0.250) {$u_j$};
        \node[font=\scriptsize] at (1.500,-0.750) {$j\in N$};
    \end{tikzpicture} \\
    \vspace{5mm}

    \begin{tikzpicture}[axpic]
        \pic (g1) at (0.500,1.000) {ge};
        \pic (g2) at (0.500,0.000) {ge};
        \pic (c1) at (1.750,0.500) {cocopy};
        \pic (c2) at (3.000,0.500) {copy};
        \pic (h1) at (4.250,1.000) {ge};
        \pic (h2) at (4.250,0.000) {ge};
        \draw (g1-out-0) to[out=0,in=180] (c1-in-0);
        \draw (g2-out-0) to[out=0,in=180] (c1-in-1);
        \draw (c1-out-0) -- (c2-in-0);
        \draw (c2-out-0) to[out=0,in=180] (h1-in-0);
        \draw (c2-out-1) to[out=0,in=180] (h2-in-0);
        \node[font=\scriptsize] at (0.050,1.200) {$x_1$};
        \node[font=\scriptsize] at (0.050,-0.250) {$m_j$};
        \node[font=\scriptsize] at (4.950,1.200) {$m_i$};
        \node[font=\scriptsize] at (4.950,-0.250) {$m_{i'}$};
    \end{tikzpicture}
    \;\begin{tikzpicture}[axpic]\node{$\overset{(\text{spider})}{=}$};\end{tikzpicture}\;
    \begin{tikzpicture}[axpic]
        \pic (c1) at (0.500,1.500) {copy};
        \pic (c2) at (0.500,-0.500) {copy};
        \pic (g11) at (2.000,2.000) {ge};
        \pic (g12) at (2.000,1.000) {ge};
        \pic (g21) at (2.000,0.000) {ge};
        \pic (g22) at (2.000,-1.000) {ge};
        \pic (d1) at (3.750,1.500) {cocopy};
        \pic (d2) at (3.750,-0.500) {cocopy};
        \draw (c1-out-0) to[out=0,in=180] (g11-in-0);
        \draw (c1-out-1) to[out=0,in=180] (g12-in-0);
        \draw (c2-out-0) to[out=0,in=180] (g21-in-0);
        \draw (c2-out-1) to[out=0,in=180] (g22-in-0);
        \draw (g11-out-0) to[out=0,in=180] (d1-in-0);
        \draw (g21-out-0) to[out=0,in=180] (d1-in-1);
        \draw (g12-out-0) to[out=0,in=180] (d2-in-0);
        \draw (g22-out-0) to[out=0,in=180] (d2-in-1);
        \node[font=\scriptsize] at (0.050,1.700) {$x_1$};
        \node[font=\scriptsize] at (0.050,-0.750) {$m_j$};
        \node[font=\scriptsize] at (4.450,1.700) {$m_i$};
        \node[font=\scriptsize] at (4.450,-0.750) {$m_{i'}$};
    \end{tikzpicture} \\
    \vspace{5mm}

    \begin{tikzpicture}[axpic]
        \pic (g1) at (0.500,0.000) {antipode};
        \pic (g2) at (2.000,0.500) {multiplication};
        \pic (g3) at (3.250,0.500) {ge};
        \pic (g4) at (4.500,0.500) {cozero};
        \coordinate (id_in_1)  at (0.000,1.000);
        \coordinate (id_out_1) at (1.250,1.000);
        \draw (id_in_1) -- (id_out_1);
        \draw (id_out_1) to[out=0,in=180] (g2-in-0);
        \draw (g1-out-0) to[out=0,in=180] (g2-in-1);
        \draw (g2-out-0) to[out=0,in=180] (g3-in-0);
        \draw (g3-out-0) to[out=0,in=180] (g4-in-0);
        \node[font=\scriptsize] at (0.050,1.200) {$x_1$};
        \node[font=\scriptsize] at (0.050,-0.250) {$m_i$};
    \end{tikzpicture}
    \;\begin{tikzpicture}[axpic]\node{$\overset{(\text{G1})}{=}$};\end{tikzpicture}\;
    \begin{tikzpicture}[axpic]
        \pic[val=\widetilde{M}] (g1) at (1.000,0.000) {2dmatrix};
        \pic (g2) at (2.400,0.000) {ge};
        \pic (g3) at (3.650,0.000) {cozero};
        \coordinate (id_in_1) at (0.000,0.300);
        \coordinate (id_in_2) at (0.000,-0.300);
        \draw (id_in_1) -- (g1-in-0);
        \draw (id_in_2) -- (g1-in-1);
        \draw (g1-out-0) to[out=0,in=180] (g2-in-0);
        \draw (g2-out-0) to[out=0,in=180] (g3-in-0);
    \end{tikzpicture}
    \;\begin{tikzpicture}[axpic]\node{$\overset{(\text{cone form})}{=}$};\end{tikzpicture}\;
    \begin{tikzpicture}[axpic]
        \pic[val=M'] (g1) at (1.000,0.000) {2dmatrix};
        \pic (g2) at (2.400,0.000) {ge};
        \pic (g3) at (3.650,0.000) {cozero};
        \coordinate (id_in_1) at (0.000,0.300);
        \coordinate (id_in_2) at (0.000,-0.300);
        \draw (id_in_1) -- (g1-in-0);
        \draw (id_in_2) -- (g1-in-1);
        \draw (g1-out-0) to[out=0,in=180] (g2-in-0);
        \draw (g2-out-0) to[out=0,in=180] (g3-in-0);
    \end{tikzpicture}
\end{center}
